\documentclass[11pt,a4paper]{article}
\usepackage[margin=27mm]{geometry}
\usepackage{amsmath,amssymb,amsthm,booktabs,longtable,hyperref}
\newtheorem{theorem}{Theorem}
\newtheorem{lemma}[theorem]{Lemma}
\newtheorem{corollary}[theorem]{Corollary}
\newcommand{\F}{\mathbb F}
\newcommand{\Sz}{\operatorname{Sz}}
\newcommand{\gen}[1]{\langle #1\rangle}
\newcommand{\Aut}{\operatorname{Aut}}
\newcommand{\N}{\operatorname{N}}
\newcommand{\C}{\operatorname{C}}
\title{Normalizer fibres in rank-two Suzuki geometries}
\author{Human Compression project}\date{}
\begin{document}\maketitle
\begin{abstract}
The families with a dihedral or a nonsplit torus normalizer as one
parabolic can be recovered from two elementary constructions. A
normalizer calculation gives uniqueness of the relevant subfield
overgroups. The remaining configurations are fibres of an incidence
relation between conjugacy classes of subgroups; these fibres are small
quotients of a Suzuki root group. Their field automorphisms and inversion
give the counts in the prime cases. This treats source Lemmas 6.2--6.6
and 6.8, with a correction to the prime-divisor condition in Table 6.
The classical subgroup classification quoted as source Lemma 6.7 is
retained as an input. The families with a subfield group as the only
maximal parabolic, and the global classification, are not concluded here.
\end{abstract}

\section{Scope and structural inputs}
Write $q=2^n$, where $n\ge3$ is odd, and put
\[
 \theta=2^{(n+1)/2},\qquad a_\epsilon(q)=q+\epsilon\theta+1
 \quad(\epsilon\in\{+1,-1\}).
\]
Let $G=\Sz(q)$. A configuration in this note is an unordered pair
of proper subgroups $\{M,L\}$, with intersection $J=M\cap L$.
Equivalence means the action of $\Aut(G)$, allowing interchange of
the two subgroups. This is the specified-group convention of the
source \cite{source}; we do not identify geometries by arbitrary
permutations which need not preserve their specified $G$ action.

The following classical structure is retained, with its attribution
to Suzuki, Tits and L\"uneburg as in Sections 2.2 and 4 of the
source. The nonidentity elements of $G$ are partitioned by its Sylow
two-subgroups and its cyclic tori of orders $q-1,a_+(q),a_-(q)$.
The normalizers of these parts are respectively
\[
 B=P\rtimes\F_q^*,\quad D=C_{q-1}\rtimes C_2,
 \quad N_\epsilon=C_{a_\epsilon(q)}\rtimes C_4.
                                                        \tag{1}
\]
These groups are maximal in $G$. Each indicated type of part forms
one conjugacy class. Every subgroup of $G$ either lies in one of
these normalizers or is a conjugate of a standard subfield group
$S=\Sz(s)$, where $s=2^d\ge8$ and $d\mid n$. In the latter case
the conjugacy class for each $s$ is unique. The same description
holds inside $S$; its maximal proper simple subfield groups have
prime field index. This is precisely the classical subgroup input
quoted from Suzuki \cite{suzuki} as source Lemma 6.7, together with
the maximal-subgroup list in Section 4. It is not proved anew here.

We also retain the standard action on the Suzuki ovoid: $B$ fixes
one point, a nonidentity involution fixes just one point, and a split
torus fixes two points whose stabilizers are interchanged by $D$.
The centralizer of an involution is its Sylow two-subgroup. Finally
$\Aut(G)=G\rtimes\gen{\phi}$, where $\phi$ is coordinate Frobenius
of order $n$.

We use the previously reviewed preliminary note for the pairwise
coprimality of the three odd orders, their proper-subfield arithmetic,
and the maximal subgroups of the solvable groups in (1). Explicitly,
in $C_t\rtimes C_4$ these maximal subgroups are
$C_t\rtimes C_2$ and $C_{t/p}\rtimes C_4$, with $p\mid t$ prime;
in $D$ they are $C_{q-1}$ and $D_{2(q-1)/p}$. In a point
stabilizer they are $P\rtimes C_{(q-1)/p}$ and $Z\rtimes C_{q-1}$,
where $Z=Z(P)$. We write $D_{2m}$ for a dihedral group of order $2m$.
The arithmetic also shows that if any one of $q-1,a_+(q),a_-(q)$
is prime, then $n$ is prime: if $n$ were composite, a proper
subfield of order at least eight would make all three composite.

For rank two, the elementary coset criterion we use is as follows.
If $M$ is maximal in $G$, $L\nleq M$, and $J=M\cap L$ is maximal
in both $M$ and $L$, then $\Gamma(G;M,L)$ is firm, connected,
flag transitive and residually weakly primitive. Indeed, the local
degrees are $[M:J]$ and $[L:J]$, both at least two; rank-two coset
geometries are flag transitive; $\gen{M,L}=G$ gives connectivity;
and the three required primitive coset actions are those with
stabilizers $M<G$, $J<M$, and $J<L$. Conversely, these maximality
conditions on $J$ are necessary when $M$ is a maximal parabolic.

\section{Two root-group quotients}
Use the standard root coordinates
\[
 P=\{u(a,b):a,b\in F\},\qquad F=\F_q,
 \qquad u(a,b)u(c,e)=u(a+c,b+e+ac^\theta).
                                                        \tag{2}
\]
Put $z=u(0,1)$, $w=u(1,0)$, $J_2=\gen z$ and $J_4=\gen w$.
Then $w^2=z$ and both $J_2,J_4$ are fixed pointwise by $\phi$.
We retain the classical realization of (2) in the faithful Suzuki
matrices, rather than reconstructing the ambient simple group.

\begin{lemma}\label{root}
The normalizers of these two subgroups are
\[
 \N_G(J_2)=P,\qquad
 \N_G(J_4)=R:=\{u(a,b):a\in\F_2,\ b\in F\}.
                                                        \tag{3}
\]
The group $R$ is abelian of order $2q$ and is isomorphic to
$C_4\times C_2^{n-1}$. The quotient $R/J_4$ is naturally the
additive group $F/\F_2$. The cyclic subgroups of order four
containing $J_2$ are parametrized by $F/\F_2$, the four-groups
containing $J_2$ by $(F/\F_2)\setminus\{0\}$, and the subgroups
of order eight containing $J_4$ by the same nonzero quotient.
These parametrizations commute with Frobenius. Also
\[
       P^\phi=R^\phi=J_4.
                                                        \tag{4}
\]
\end{lemma}
\begin{proof}
The first equality in (3) is the retained involution-centralizer
fact, since an automorphism of a group of order two is trivial.
Every element normalizing $J_4$ fixes $w^2=z$, so lies in $P$.
The multiplication law gives
\[
 u(a,b)^2=u(0,a^{\theta+1}),\qquad
 [u(a,b),w]=u(0,a+a^\theta),
                                                        \tag{5}
\]
where the choice of commutator convention does not change the
second formula because its value has order at most two.
Conjugation preserves the first coordinate of $w$, so normalizing
$J_4$ is equivalent to $a+a^\theta\in\F_2$.
The absolute trace of $a+a^\theta$ is zero, whereas the trace of
$1$ is one because $n$ is odd. Thus $a=a^\theta$. The fixed
field of $x\mapsto x^\theta$ is $\F_2$, since
$\gcd(n,(n+1)/2)=1$. This proves (3). Formula (2) on first
coordinates in $\F_2$ proves that $R$ is abelian, has exponent
four, and has squares exactly $J_2$. Its stated abstract structure
and $R/J_4\cong F/\F_2$ follow, the latter map being represented
by $u(0,b)\mapsto b+\F_2$.

The exponent $\theta+1$ is relatively prime to $q-1$: a common
odd divisor would give $\theta=-1$ and $\theta^2=2$ modulo that
divisor, hence $1=2$. Thus (5) shows that an element whose square
is $z$ has first coordinate one. The groups $\gen{u(1,b)}$ are
therefore all cyclic groups of order four in $P$ containing $J_2$, and
two values of $b$ give the same group precisely when they differ
by one. A cyclic group of order four containing $J_2$ centralizes
its square $z$, so is in $P$. Likewise, a four-group containing
$J_2$ centralizes $z$ and lies in $P$. All involutions of $P$
belong to $Z=\{u(0,b):b\in F\}$, so
four-groups containing $J_2$ are the two-dimensional binary
subspaces of $F$ containing one. These are parametrized by the
nonzero elements of $F/\F_2$.

If a group of order eight contains $J_4$, that index-two subgroup
is normal; the group is therefore in $R$. Conversely, inverse
images of order-two subgroups of $R/J_4$ have order eight.
They are $C_4\times C_2$ since $R$ is abelian of exponent four.
All formulas commute with squaring of field coordinates. Formula
(4) follows because the only fixed coordinates are zero and one.
\end{proof}

\begin{lemma}\label{fixed}
Suppose $n$ is prime. Frobenius has the following numbers of fixed
cosets:
\[
\begin{array}{c|ccc}
 X&P/J_2&P/J_4&R/J_4\\\hline
 |X|&q^2/2&q^2/4&q/2\\
 |X^\phi|&2&1&1.
\end{array}                                               \tag{6}
\]
On $(F/\F_2)\setminus\{0\}$ it has no fixed point. On
$(P/J_2)\setminus\{1\}$ modulo inversion, it has exactly one
fixed orbit. The number of inversion orbits in this last set is
\[
                         q^2/4+q/2-1.                  \tag{7}
\]
\end{lemma}
\begin{proof}
For either pointwise fixed subgroup $J=J_2,J_4$, a fixed right
coset $gJ$ gives $\phi(g)=gh$ for some $h\in J$.
Iterating $n$ times gives $h^n=1$, so $h=1$ because $n$ is odd
and $h$ has two-power order. Hence fixed cosets are exactly the
cosets represented by fixed elements. Equation (4) proves (6).
The statement for $F/\F_2$ follows from $R/J_4$.

Let $H=P/J_2$. It has $q$ elements with square one: by (5),
there are $2q$ elements of $P$ with square in $J_2$, with first
coordinate zero or one, and each coset has two elements.
Inversion on the set $H$ fixes exactly these $q$ elements.
Pairing the others and removing the identity gives
$(|H|+q)/2-1$, which is (7). An inversion orbit fixed by
$\phi$ has size at most two. The permutation of it induced by
an element of odd prime order cannot exchange its two members;
each member is therefore fixed. By (6), $H^\phi$ has order two,
and its sole nonidentity element gives the one fixed orbit.
\end{proof}

\section{Uniqueness of three subfield overgroups}
\begin{lemma}\label{incidence}
Let a finite group $K$ have one conjugacy class of subgroups $S$
of the specified type. Suppose $H<S$, and all the $K$-conjugates
of $H$ lying in $S$ form one $S$-conjugacy class. Then the number of
conjugates of $S$ containing a fixed conjugate of $H$ is
\[
 \frac{|S|\,|\N_K(H)|}{|\N_K(S)|\,|\N_S(H)|}.
                                                        \tag{8}
\]
In particular it is one if $S$ and $H$ are self-normalizing in
$K$.
\end{lemma}
\begin{proof}
Count incident pairs of conjugates $(S',H')$ with $H'\le S'$.
There are $|K|/|\N_K(S)|$ first entries, each incident with
$|S|/|\N_S(H)|$ second entries. Divide by the total number
$|K|/|\N_K(H)|$ of second entries.
\end{proof}

\begin{lemma}\label{embeddings}
For each subfield order $s\ge8$, every subgroup abstractly
isomorphic to one of the three groups in (9) below is conjugate
to its standard copy in $\Sz(s)$.
\end{lemma}
\begin{proof}
For a dihedral group or a group $C_{a_\delta(s)}\rtimes C_4$,
the odd cyclic subgroup is characteristic. Its nonidentity
elements determine the unique ambient torus. The order of that
torus is the same as for the standard copy: the three ambient
odd orders are pairwise coprime. All tori of a given type are
conjugate, and each has a unique subgroup of each dividing
order. The complements are conjugate by the torus. In the
dihedral case this is the invertibility of squaring on an odd
cyclic group; in the order-four case it follows from the
invertibility of $q-1$ on the nonsplit torus. This proves the
assertion for the two cyclic normalizers.

Now let $H$ be abstractly isomorphic to $B_s$. Its normal Sylow
two-subgroup $V$ has order $s^2$, with exactly $s-1$ involutions,
and its odd complement $A$ has order $s-1$. The centre of $V$
contains an involution $z'$, so $V\le\C_G(z')$, a Sylow
two-subgroup. Any nonidentity element of $V$ is in a unique
part of the classical partition. Consequently $A$ normalizes
this Sylow subgroup, and after conjugation $H\le B=P\rtimes T$,
where $T$ denotes the diagonal torus.
The image of $A$ in $B/P=T$ is its unique subgroup $K^*$,
where $K=\F_s$. Conjugating by an element of $P$ makes
$A=K^*$: first apply the affine complement lemma to
$(P/Z)\rtimes K^*$, and then to $Z\rtimes K^*$. The lemma
permits a nonfaithful action, so applies to both quotients.

All involutions of $P$ are in $Z$. Hence $|V\cap Z|=s$ and
the image of $V$ in $P/Z$ also has order $s$. The action of
$K^*$ on each coordinate is the full scalar group of $K$.
Thus these two additive subgroups are lines $bK$ and $aK$
in $F$. Squares in (5) give $bK=a^{\theta+1}K$. A diagonal
conjugation makes $aK=K$ and consequently $bK=K$.
The possible subgroup $V$ is now described by a function
$f:K\longrightarrow F/K$:
\[
       V=\{u(c,e):c\in K,\ e+K=f(c)\}.
\]
The multiplication law shows that $f$ is additive, since
$c d^\theta\in K$ for $c,d\in K$. Invariance under $K^*$
gives $f(k^\theta c)=k^{\theta+2}f(c)$. Since
$\theta^2=2$ as an exponent on $K^*$, this implies
$f(c)=c^{\theta+1}f(1)$ for every $c\in K$.
Choose $c\in K\setminus\F_2$. Additivity at $c+1$ gives
$(c^\theta+c)f(1)=0$ in the $K$-vector space $F/K$.
The coefficient is nonzero, because the fixed field of
$x\mapsto x^\theta$ on $K$ is $\F_2$. Therefore $f(1)=0$,
so $f=0$ and $V=U(K,K)$. This proves $H=B_s$ after conjugation.
\end{proof}

\begin{theorem}\label{unique}
Let $K=\F_s\subseteq F$, with $s\ge8$. Each subgroup of one
of the following types is contained in exactly one subgroup
isomorphic to $\Sz(s)$:
\[
 B_s=U(K,K)\rtimes K^*,\qquad
 D_{2(s-1)},\qquad C_{a_\delta(s)}\rtimes C_4
 \quad(\delta=+1,-1).
                                                        \tag{9}
\]
\end{theorem}
\begin{proof}
By Lemma~\ref{embeddings}, it suffices to use the groups in (9)
inside the standard $S=\Sz(s)$.
Their respective conjugacy classes in $S$ are unique by the
retained structure. We prove that each is self-normalizing in
$G$. The subfield group $S$ is also self-normalizing. Indeed,
its normalizer contains the nonsolvable simple group $S$, so
cannot be contained in a solvable normalizer in (1). The subgroup
classification makes it a simple subfield group; since it
normalizes its nontrivial normal subgroup $S$, it equals $S$.

For $B_s$, its characteristic Sylow two-subgroup determines the
unique fixed ovoid point, so any normalizer lies in $B=PH$.
Write its elements using $u(a,b)$ and $h(t)$, where
\[
 u(a,b)^{h(t)}=u(t^\theta a,t^{\theta+2}b).
                                                        \tag{10}
\]
The exponent $\theta+2$ is coprime to $q-1$: a common odd
divisor would give $\theta=-2$ and $\theta^2=2$, hence $4=2$.
Preserving the centre $U(0,K)$ of the Sylow subgroup in $B_s$
therefore forces $t^{\theta+2}\in K^*$, and hence $t\in K^*$.
Remove this factor. For $k\in K^*\setminus\{1\}$, the first
coordinate of the commutator of $u(a,b)$ with $h(k)$ is a
nonzero $K$-multiple of $a$. Membership in $B_s$ forces $a\in K$.
Remove the factor $u(a,0)\in B_s$ as well. The remaining central
factor $u(0,b)$ has commutator coordinate
$(k^{\theta+2}-1)b$, up to multiplication by a nonzero element
of $K$. This forces $b\in K$. Thus $\N_G(B_s)=B_s$.

For either odd normalizer in (9), its cyclic odd subgroup is
characteristic. The partition into tori forces its normalizer
to preserve the unique ambient torus containing that subgroup.
In the split case this gives a subgroup of $C_{q-1}\rtimes C_2$.
A torus element $x$ normalizes $D_{2(s-1)}$ precisely when
$x^2\in C_{s-1}$; since the torus has odd order this is equivalent
to $x\in C_{s-1}$.

In the nonsplit case the ambient torus normalizer is $C_t\rtimes
\gen v$, with $t=a_\epsilon(q)$ and $|v|=4$. The standard
action on the cyclic torus is the power $q$ or its inverse.
Both $q-1$ and $q^{-1}-1$ are coprime to $t$, because
$\gcd(q-1,t)=1$. For $m=a_\delta(s)$, a torus element $x$
normalizes $C_m\rtimes\gen v$ precisely when the commutator
of $x$ with $v$ lies in $C_m$. This is equivalent to $x\in C_m$.
All complements of order four are torus conjugate, so this
calculation applies to the subgroup in (9). It is self-normalizing.
Equation (8) now proves all three uniqueness assertions.
\end{proof}

These are source Lemmas 6.5, 6.6 and 6.8, including the version
of 6.8 for the other nonsplit torus. The common step is incidence
counting after self-normalization. No new conjugacy classification
of subfield groups is being asserted.

\section{Fibres of subgroup incidence}
Fix $J=J_2$ or $J_4$. If the copies of $J$ in a subgroup $M$ form
one conjugacy class in $M$, then the conjugates of $M$ containing
this fixed $J$ form one $\N_G(J)$-orbit. To prove this, first
conjugate one copy of $M$ to another; conjugate again within the
second copy to send the image of $J$ to $J$. The stabilizer is
$\N_G(J)\cap\N_G(M)$. Thus the orbit is an actual coset space,
not merely a set of the same cardinality.

Applied to (1) and to its prime subgroups, this gives
\[
\begin{array}{c|c|c|c}
 J&M&\N_G(J)\cap\N_G(M)&\text{fibre}\\\hline
 J_2&D_{2(q-1)}&J_2&P/J_2\\
 J_2&D_{2p},\ p\mid a_\epsilon(q)\text{ prime}&C_4&P/C_4\\
 J_4&C_t\rtimes C_4&J_4&R/J_4.
\end{array}                                               \tag{11}
\]
In the last row $t$ may be the whole nonsplit torus order or a
prime divisor of it. For the second row, the full normalizer of
$D_{2p}$ is $C_p\rtimes C_4$. This follows by preserving its
characteristic odd subgroup, then using the odd-order squaring
argument from Theorem~\ref{unique}. Its stabilizer of its given
involution is a cyclic group of order four. For the last row,
the commutator argument in the same theorem shows that
$C_t\rtimes C_4$ is self-normalizing, even when $t$ is a proper
prime divisor of the ambient torus order. The required conjugacy
of involutions or complements inside these subgroups follows
from the cyclic semidirect products of odd order by $C_2$ or
$C_4$. This also follows from the affine complement argument
in the preliminary note.

Suppose now that $n$ is prime. Every fibre in (11) has a
Frobenius-fixed member: its cardinality is respectively $q^2/2$,
$q^2/4$ or $q/2$, congruent to $2,1,1$ modulo $n$ by Fermat's
congruence. The orbits of a cyclic group of prime order have size
one or $n$, so a fixed member exists. Choose that member as base.
In the second row its stabilizer $C_4$ is Frobenius invariant,
and Frobenius acts trivially on it because $\Aut(C_4)$ has order
two. By (4), this stabilizer is the particular $J_4$ used above.
Thus (11) is Frobenius equivariant with exactly the quotients
in (6); no assumption about a preferred complement is needed.

For two distinct dihedral parabolics with common $J_2$, the first
fibre in (11) is a regular set for $H=P/J_2$. Simultaneous
translation sends an ordered pair to $(1,h)$ with $h\ne1$;
interchanging its entries sends $h$ to $h^{-1}$. Hence the
unordered configurations are the inversion classes in
$H\setminus\{1\}$, followed by Frobenius. For two equal types
of nonsplit normalizer, the same argument uses $R/J_4$, an
elementary abelian group, so inversion does nothing.

For distinct parabolic types there is no interchange. Once the
first parabolic containing $J_2$ is fixed, its stabilizer in $P$
is $J_2$, which acts trivially on $P/J_4$ and on the order-four
groups in Lemma~\ref{root}. The nonsplit calculation has the
analogous trivial action of $J_4$ on $R/J_4$. These observations
explain why the small fibres give the complete configuration
counts, rather than just numbers of possible second subgroups.

\section{The split normalizer}
\begin{theorem}\label{split}
Take $M=D_{2(q-1)}$. If $q-1$ is prime, the possible second
parabolics and their numbers of equivalence classes are precisely
the rows of the following table. Each prime $p$ in the third row
gives a separate family.
\[
\begin{array}{c|c|c}
 L&J&\text{number}\\\hline
 Z\rtimes C_{q-1}&C_{q-1}&1\\
 D_{2(q-1)}&C_2&1+\dfrac{(q-2)(q+4)}{4n}\\[2mm]
 D_{2p},\ p\mid a_+(q)a_-(q)&C_2&1+\dfrac{(q-2)(q+2)}{4n}\\[2mm]
 C_2^2&C_2&\dfrac{q-2}{2n}\\[2mm]
 C_4&C_2&1+\dfrac{q-2}{2n}.
\end{array}                                               \tag{12}
\]
If $q-1$ is composite, the first row still occurs once. The only
other possibilities are
\[
 L=\Sz(s),\qquad J=D_{2(s-1)},
 \qquad (q-1)/(s-1)\text{ prime},                         \tag{13}
\]
with $s\ge8$ a proper subfield order. Each occurs once.
\end{theorem}
\begin{proof}
First consider $J=C_{q-1}=A$, which is always maximal in $M$.
A proper overgroup $L$ of $A$ not in $M$ must lie in a point
stabilizer, by the subgroup classification and coprimality of
the three odd factors. A proper subfield group cannot contain
$A$: if $q\ge s^3$, every cyclic torus in $\Sz(s)$ has order
less than $q-1$. The torus $A$ fixes two ovoid points, giving
the two possible point stabilizers.

In either one write $L=VA$, with $V=L\cap P$ invariant under
the full scalar torus. If $V\le Z$, nonzero invariance forces
$V=Z$, since the action is all nonzero scalars. If $V\nleq Z$,
the same argument on $P/Z$ shows $VZ=P$. The Frattini property
$Z=\Phi(P)$ from the preliminary note then gives $V=P$.
But $A<ZA<PA$, so $A$ would not be maximal in $L$.
Consequently $L=ZA$. The two choices are interchanged by $M$,
and their intersection with $M$ is exactly $A$. This proves the
first row for every $q$.

Suppose $q-1$ is prime. The other maximal intersection in $M$
is $J_2$. We classify proper overgroups $L$ in which $J_2$ is
maximal. If $L$ is contained in a point stabilizer, its normal
two-subgroup $L\cap P$ contains $J_2$. By maximality it either
equals $J_2$ or equals $L$. In the first case $L$ normalizes
$J_2$, hence lies in its two-group centralizer $P$ as well.
In both cases $L$ is a two-group, and maximality makes its
order four. It is therefore $C_4$ or $C_2^2$.
The cyclic normalizers give $D_{2p}$, with $p$ prime, or $C_4$.
The split case has $p=q-1$; the other two odd orders give the
third row. No simple subfield group has a maximal subgroup of
order two, by its retained maximal-subgroup list. This proves
exhaustion. The intersections in all distinct-type cases are
exactly $J_2$ by their orders. Two distinct split normalizers
containing $J_2$ cannot have a common odd element, since such
an element lies in a unique torus; their intersection is also
exactly $J_2$. All the groups are outside $M$, so the coset
criterion applies.

Here $n$ is prime. The inversion quotient for the second row
has size (7) and one fixed point by Lemma~\ref{fixed}. An action
of order $n$ on a finite set $X$ has
$(|X|+(n-1)|X^\phi|)/n$ orbits. This gives
$1+(q-2)(q+4)/(4n)$. The third row uses $P/J_4$, of size
$q^2/4$ with one fixed point, giving
$1+(q-2)(q+2)/(4n)$. The last two rows use respectively the
nonzero quotient $F/\F_2$ and the whole quotient, of size
$q/2-1$ and $q/2$, with zero and one fixed points. This proves
all counts in (12).

Finally suppose $q-1$ is composite. Apart from $A$, a maximal
intersection has form $J=D_{2m}$, with $m=(q-1)/p>1$ and $p$
prime. Any overgroup of $J$ lying in a solvable normalizer must
lie in $M$: the characteristic odd rotation subgroup determines
its ambient split torus; a point stabilizer cannot contain its
inverting involution, since its quotient by the normal two-group
is abelian of odd order. Coprimality excludes nonsplit tori.
Thus an overgroup outside $M$ is $L=\Sz(s)$. Maximality of
$J$ in $L$ forces $m=s-1$ by its classical list. Conversely
the standard such $J$ occurs when the condition in (13) holds.
It lies in a unique $L$ by Theorem~\ref{unique}. Its conjugacy
class in $M$ is unique, so this produces exactly one equivalence
class. The maximality conditions and exact intersection follow
as before.
\end{proof}

Theorem~\ref{split} recovers source Lemma 6.2 and Tables 2 and 3.
The common proof does not separately enumerate the subgroups of
each ambient Suzuki group.

\section{The nonsplit normalizers}
\begin{theorem}\label{nonsplit}
Let $t=a_\epsilon(q)$ and $M=C_t\rtimes C_4$. If $t$ is prime,
the complete list is
\[
\begin{array}{c|c|c}
 L&J&\text{number}\\\hline
 C_t\rtimes C_4&C_4&\dfrac{q-2}{2n}\\[2mm]
 C_p\rtimes C_4,\ p\mid a_{-\epsilon}(q)\text{ prime}
      &C_4&1+\dfrac{q-2}{2n}\\[2mm]
 C_4\times C_2&C_4&\dfrac{q-2}{2n}.
\end{array}                                               \tag{14}
\]
If $t$ is composite, the only possibilities are
\[
 L=\Sz(s),\qquad J=C_{a_\delta(s)}\rtimes C_4,
 \qquad t/a_\delta(s)\text{ prime},                       \tag{15}
\]
where $s\ge8$ is a proper subfield order and
$a_\delta(s)\mid t$. Each has one equivalence class.
\end{theorem}
\begin{proof}
One maximal subgroup of $M$ is $C_t\rtimes C_2$. Its nontrivial
characteristic odd subgroup determines the ambient nonsplit
torus. Any overgroup outside $M$ must therefore be a simple
subfield group, by the same partition and coprimality argument
as above. For $C_t\rtimes C_2$ to be maximal in $\Sz(s)$ the
classical list would require $t=s-1$. This is impossible:
$s-1\mid q-1$, whereas $\gcd(t,q-1)=1$ and $t>1$.

If $t$ is prime, the remaining maximal intersection is $J_4$.
A proper overgroup $L$ in which $J_4$ is maximal, if contained
in a point stabilizer, must be a two-group. Indeed its normal
two-subgroup contains $J_4$, and if it equals $J_4$ the whole
group lies in $\N_G(J_4)=R$, still a two-group. Maximality
makes $|L|=8$, whence Lemma~\ref{root} gives $C_4\times C_2$.
A split torus normalizer has no element of order four. A
nonsplit one gives $C_p\rtimes C_4$, with $p$ prime. The prime
$p=t$ gives another conjugate of $M$; all other primes must
divide the other nonsplit torus order. The classical maximal
list excludes a simple subfield group with maximal subgroup
$C_4$. This proves exhaustion. Distinct conjugates of $M$
containing $J_4$ intersect exactly in $J_4$, since any common
odd element would determine the same torus and normalizer.
Different types have intersection $J_4$ by their orders.

Since $n$ is prime, the first row is the nonzero elementary
abelian quotient $R/J_4$ modulo Frobenius, giving
$(q-2)/(2n)$. The second is the whole quotient, with one fixed
point, giving $1+(q-2)/(2n)$. The third is the nonzero quotient
from Lemma~\ref{root}. These are actual incidence fibres by
(11), and the interchange argument has already been included.

If $t$ is composite, the remaining maximal intersections are
$J=C_m\rtimes C_4$, where $m=t/p>1$. Their nontrivial odd
subgroups again exclude any solvable overgroup outside $M$.
Thus $L=\Sz(s)$, and maximality in its classical list forces
$m=a_\delta(s)$. Conversely each condition (15) gives the
standard subgroup and its unique subfield overgroup by
Theorem~\ref{unique}. Conjugacy within $M$ and this uniqueness
give one class. The coset criterion proves the claimed geometric
properties in every listed case.
\end{proof}

For $\epsilon=+1$, this gives source Lemma 6.3 and Tables 4--5.
For $\epsilon=-1$, it gives Lemma 6.4 and Tables 6--7, subject
to the following explicit correction. The second row of printed
Table 6, page 110, requires its prime $p$ to divide $\beta_q$,
the source's notation for $a_-(q)$. That table already assumes
$\beta_q$ prime, so this repeats the same-normalizer family in
its first row and omits the opposite torus family. The condition
must be $p\mid\alpha_q=a_+(q)$, as (14) proves. At $q=8$ the
missing family has parabolics $C_5\rtimes C_4$ and
$C_{13}\rtimes C_4$. It is not a new conjecture or a global
classification correction; it is a corrected condition in this
particular table.

\section{Attribution and remaining obligations}
The recovered assertions are source Lemmas 6.2--6.6 and 6.8,
and Tables 2--7, with the Table 6 condition just specified.
The source's Lemma 6.7 remains an explicitly attributed classical
input from Suzuki. Our contribution in this component is the
common derivation through normalizer incidence, the two root-group
quotients and their Frobenius actions. We make no priority claim
for incidence counting, Burnside's orbit formula, or these elementary
group-theoretic tools. The source's original case proofs and numerical
assertions remain attributed to \cite{source}.

No part of this note supplies the remaining global classification.
In particular, configurations whose only maximal parabolic is a
subfield group require separate treatment of all its maximal
intersections. The already reviewed scalar-subspace correction
addresses one such family and shows that the printed global
classification cannot be retained unchanged. Neither the present
note nor that correction alone settles all remaining cases or
the source's global numerical tables. The attached finite controls
at $q=8,32,128$ form the literal root groups, their indicated
subgroups and coset partitions, and the Frobenius and inversion
orbits. All seven count formulas agree. These checks concern the
quotient algebra and do not test the borrowed classification of
ambient subgroups or replace the quantified arguments. This component
has no Lean certificate and carries no whole-original completion
credit by itself.

\begin{thebibliography}{9}
\bibitem{source} D.~Leemans, \emph{The rank 2 geometries of the simple
Suzuki groups $\Sz(q)$}, Beitr\"age zur Algebra und Geometrie
39 (1998), no.~1, 97--120.
\bibitem{suzuki} M.~Suzuki, \emph{On a class of doubly transitive groups},
Annals of Mathematics 75 (1962), 105--145.
\end{thebibliography}
\end{document}
